
\documentclass[a4paper, 10pt, twoside]{article}

\usepackage[top=1in, bottom=1in, left=1in, right=1in]{geometry}
\usepackage[utf8]{inputenc}
\usepackage[spanish, es-ucroman, es-noquoting]{babel}
\usepackage{setspace}
\usepackage{fancyhdr}
\usepackage{lastpage}
\usepackage{amsmath}
\usepackage{amsfonts}
\usepackage{amsthm}
\usepackage{verbatim}
\usepackage{fancyvrb}
\usepackage{graphicx}
\usepackage{float}
\usepackage{enumitem} % Provee macro \setlist
\usepackage{tabularx}
\usepackage{multirow}
\usepackage{hyperref}
\usepackage{xspace}
\usepackage{qtree}
\usepackage[toc, page]{appendix}


%%%%%%%%%% Constantes - Inicio %%%%%%%%%%
\newcommand{\titulo}{Trabajo Práctico}
\newcommand{\materia}{Teoría de Lenguajes}
\newcommand{\integrantes}{Delgado · Lovisolo · Petaccio}
\newcommand{\cuatrimestre}{Segundo Cuatrimestre de 2014}
%%%%%%%%%% Constantes - Fin %%%%%%%%%%


%%%%%%%%%% Configuración de Fancyhdr - Inicio %%%%%%%%%%
\pagestyle{fancy}
\thispagestyle{fancy}
\lhead{\titulo\ · \materia}
\rhead{\integrantes}
\renewcommand{\footrulewidth}{0.4pt}
\cfoot{\thepage /\pageref{LastPage}}

\fancypagestyle{caratula} {
   \fancyhf{}
   \cfoot{\thepage /\pageref{LastPage}}
   \renewcommand{\headrulewidth}{0pt}
   \renewcommand{\footrulewidth}{0pt}
}
%%%%%%%%%% Configuración de Fancyhdr - Fin %%%%%%%%%%


%%%%%%%%%% Miscelánea - Inicio %%%%%%%%%%
% Evita que el documento se estire verticalmente para ocupar el espacio vacío
% en cada página.
\raggedbottom

% Separación entre párrafos.
\setlength{\parskip}{0.5em}

% Separación entre elementos de listas.
\setlist{itemsep=0.5em}

% Asigna la traducción de la palabra 'Appendices'.
\renewcommand{\appendixtocname}{Apéndices}
\renewcommand{\appendixpagename}{Apéndices}

\newcommand{\grafico}[3]{
  \begin{figure}[H]
  	\centering
    \includegraphics[height=10cm]{#1}
    \caption{#2}
    \label{#3}
  \end{figure}
}


%%%%%%%%%% Miscelánea - Fin %%%%%%%%%%


\begin{document}


%%%%%%%%%%%%%%%%%%%%%%%%%%%%%%%%%%%%%%%%%%%%%%%%%%%%%%%%%%%%%%%%%%%%%%%%%%%%%%%
%% Carátula                                                                  %%
%%%%%%%%%%%%%%%%%%%%%%%%%%%%%%%%%%%%%%%%%%%%%%%%%%%%%%%%%%%%%%%%%%%%%%%%%%%%%%%


\thispagestyle{caratula}

\begin{center}

\includegraphics[height=2cm]{DC.png} 
\hfill
\includegraphics[height=2cm]{UBA.jpg} 

\vspace{2cm}

Departamento de Computación,\\
Facultad de Ciencias Exactas y Naturales,\\
Universidad de Buenos Aires

\vspace{4cm}

\begin{Huge}
\titulo
\end{Huge}

\vspace{0.5cm}

\begin{Large}
\materia
\end{Large}

\vspace{1cm}

\cuatrimestre

\vspace{4cm}

\begin{tabular}{|c|c|c|}
\hline
Apellido y Nombre & LU & E-mail\\
\hline
Delgado, Alejandro N.  & 601/11 & nahueldelgado@gmail.com\\
Lovisolo, Leandro      & 645/11 & leandro@leandro.me\\
Petaccio, Lautaro José & 443/11 & lausuper@gmail.com\\
\hline
\end{tabular}

\end{center}

\newpage


%%%%%%%%%%%%%%%%%%%%%%%%%%%%%%%%%%%%%%%%%%%%%%%%%%%%%%%%%%%%%%%%%%%%%%%%%%%%%%%
%% Introducción                                                              %%
%%%%%%%%%%%%%%%%%%%%%%%%%%%%%%%%%%%%%%%%%%%%%%%%%%%%%%%%%%%%%%%%%%%%%%%%%%%%%%%


\section{Introducción}
Presentaremos en este informe, un intérprete encargado de procesar el lenguaje provisto por la cátedra y generar los renders tridimensionales asociados a este lenguaje.

Utilizaremos la librería Ply, sobre la cuál definiremos la gramática para este lenguaje, para luego construir el lexer y el parser que completarán el funcionamiento del intérprete.


\section{Especificación de la gramática}

A continuación detallamos primero los tokens reconocidos por el lexer y luego las producciones reconocidas por el parser.


\subsection{Tokens}

Las cadenas representadas por cada token se corresponden con su nombre en minúscula salvo especificado lo contrario.


\vspace{1em}
\begin{tabular}{ll}
\hline
Token & Descripción\\
\hline\hline
\texttt{BOX}, \texttt{BALL}                       & Primitivas\\
\texttt{UNDERSCORE}                               & Primitiva vacía (carácter ``\_'')\\
\texttt{CYLINDER}, \texttt{CONE}, \texttt{TORUS}  & Primitivas adicionales\\
\texttt{RX}, \texttt{RY}, \texttt{RZ}             & Transformaciones de rotación\\
\texttt{SX}, \texttt{SY}, \texttt{SZ}, \texttt{S} & Transformaciones de escala\\
\texttt{TX}, \texttt{TY}, \texttt{TZ}             & Transformaciones de traslación\\
\texttt{CR}, \texttt{CG}, \texttt{CB}             & Transformaciones de coloreo\\
\texttt{D}                                        & Transformación de límite de profundidad de recursión\\
\texttt{COLON}                                    & Inicio de una transformaciín (carácter ``:'')\\
\texttt{AND}, \texttt{OR}                         & Operaciones de conjunción y disyunción (caracteres ``\&'' y ``\textbar'')\\
\texttt{POWER}                                    & Operación de potenciación (carácter ``\^{}'')\\
\texttt{LGROUP}, \texttt{RGROUP}                  & Agrupación de operaciones (caracteres ``['' y ``]'')\\
\texttt{LOPT}, \texttt{ROPT}                      & Operación opcional (caracteres ``\textless'' y ``\textgreater'')\\
\texttt{NUMBER}                                   & Constante numérica (expresión regular \texttt{[0-9]+(\textbackslash.[0-9]+)?})\\
\texttt{PLUS}, \texttt{MINUS},
\texttt{TIMES}, \texttt{DIVIDE}                   & Operadores aritméticos binarios (caracteres ``+'', ``-'', ``*'' y ``/'')\\
\texttt{LPAREN}, \texttt{RPAREN}                  & Agrupación de operaciones aritméticas (caracteres ``('' y ``)'')\\
\texttt{START\_RULE}                               & Regla inicial (carácter ``\$'')\\
\texttt{RULE}                                     & Regla (expresión regular \texttt{[a-zA-Z]+})\\
\texttt{DOT}                                      & Regla final (carácter ``.'')\\
\texttt{EQUALS}                                   & Declaración de regla (carácter ``='')\\
\texttt{COMMENT}                                  & Comentario (expresión regular \texttt{"([\^{}\textbackslash\textbackslash\textbackslash{}n]|(\textbackslash\textbackslash.))*"})\\
\hline
\end{tabular}


\subsection{Producciones}

Las producciones a continuación están escritas en el formato consumido por la librería PLY.

La producción inicial de esta gramática es \texttt{rules}.

Notar que la producción \texttt{empty} no tiene cuerpo. Esto es porque, como su nombre señala, se corresponde con la producción vacía, comúnmente escrita $A \rightarrow \epsilon$ o $A \rightarrow \lambda$ en los libros de texto, y la librería PLY no provee ninguna sintaxis especial para tales producciones.

Notar además que algunas producciones utilizan una sintaxis especial \texttt{\%proc}. Esta sintaxis se explica en la próxima sección.

\begin{verbatim}
  rules : rule_definition rules
        | empty

  rule_definition : RULE       EQUALS element
                  | RULE DOT   EQUALS element
                  | START_RULE EQUALS element

  element : primitive
          | rule
          | transform
          | element_and
          | element_or
          | element_power
          | element_group
          | element_optional

  primitive : BOX
            | BALL
            | UNDERSCORE
            | CYLINDER
            | CONE
            | TORUS

  rule : RULE

  transform : element COLON transform_name arith_expr

  transform_name : RX
                 | RY
                 | RZ
                 | SX
                 | SY
                 | SZ
                 | S
                 | TX
                 | TY
                 | TZ
                 | CR
                 | CG
                 | CB
                 | D

  element_and : element AND element

  element_or : element OR element

  element_power : element POWER arith_expr

  element_group : LGROUP element RGROUP

  element_optional : LOPT element ROPT

  arith_expr : arith_expr_number
             | arith_expr_uplus
             | arith_expr_uminus
             | arith_expr_parenthesis
             | arith_expr_plus
             | arith_expr_minus
             | arith_expr_times
             | arith_expr_divide

  arith_expr_number : NUMBER

  arith_expr_uplus : PLUS arith_expr %prec UPLUS

  arith_expr_uminus : MINUS arith_expr %prec UMINUS

  arith_expr_parenthesis : LPAREN arith_expr RPAREN

  arith_expr_plus : arith_expr PLUS arith_expr

  arith_expr_minus : arith_expr MINUS arith_expr

  arith_expr_times : arith_expr TIMES arith_expr

  arith_expr_divide : arith_expr DIVIDE arith_expr

  empty :
\end{verbatim}


\subsection{Órdenes de precedencia y asociatividades}

Las ambiguedades en las producciones anteriores se resuelven con los órdenes de precedencia y asociatividades en la tabla a continuación.

Notar que la tabla incluye dos tokens ficticios \texttt{UPLUS} y \texttt{UMINUS}. Éstos se corresponden con los operadores unarios ``+'' y ``-'', que definen el signo de una expresión aritmética. Notar que tienen un orden de precedencia superior al de sus contrapartes binarios.

Estos tokens ficticios se referencian desde las producciones \texttt{arith\_expr\_uplus} y \texttt{arith\_expr\_uminus}, que usan los tokens \texttt{PLUS} y \texttt{MINUS}, respectivamente, pero que en el contexto de esas producciones particulares requieren un orden de precedencia distinto. La librería PLY provee un mecanismo para cambiar el orden de precedencia de un token en una producción: escribiendo \texttt{\%prec \textless{}TOKEN\textgreater} al final de una producción, el orden de precedencia del token usado en esa producción se reemplaza por el orden de precedencia del token \texttt{\textless{}TOKEN\textgreater} en ese contexto.


\vspace{1em}
\begin{tabular}{llll}
  \hline
  Orden & Token           & Asociatividad & Comentarios\\
  \hline\hline
  1     & \texttt{PLUS}   & Izquierda     & Operadores aritméticos binarios\\
  1     & \texttt{MINUS}  & Izquierda     &\\
  2     & \texttt{TIMES}  & Izquierda     &\\
  2     & \texttt{DIVIDE} & Izquierda     &\\
  \hline
  3     & \texttt{UPLUS}  & Derecha       & Operadores aritméticos unarios\\
  3     & \texttt{UMINUS} & Derecha       &\\
  \hline
  4     & \texttt{AND}    & Izquierda     & Operadores binarios sobre elementos\\
  5     & \texttt{OR}     & Izquierda     &\\
  \hline
  6     & \texttt{POWER}  & Izquierda     & Operadores unarios sobre elementos\\
  7     & \texttt{COLON}  & Izquierda     &\\
  \hline
\end{tabular}


\section{Implementación}

La ejecución del programa se divide en dos fases: parsing de la entrada y rendering de la escena. Explicamos ambas a continuación.


\subsection{Parsing}

Durante esta fase se ejecuta el parser generado con la librería PLY y se arma una instancia de la clase \texttt{Scene}, que captura la semántica de la entrada del programa.

\texttt{Scene} tiene una variable de instancia \texttt{Scene::rules} que mantiene un arreglo de instancias de la clase \texttt{Rule} o alguna de sus subclases. Este arreglo tiene un item por cada regla definida en la entrada. Por ejemplo, la entrada a continuación genera un arreglo \texttt{Scene::rules} con 4 items:

\begin{verbatim}
a  = box  & a : tx 1
a  = ball & a : tx 1
a. = box : cg 0
$  = a : d 10
\end{verbatim}

Las reglas iniciales y finales corresponden a las subclases \texttt{StartRule} y \texttt{FinalRule}. El resto a la clase \texttt{Rule}.

Cada regla se modela mediante un árbol de parsing con exactamente la misma estructura reconocida por el parser. Cada nodo de este árbol se representa con una instancia de alguna subclase de \texttt{Node}, definida de la siguiente manera:

\begin{verbatim}
  class Node:
    def __init__(self, scene):
      self.scene = scene
      self.children = []
    ...
\end{verbatim}

Notar que la clase \texttt{Rule} es subclase de \texttt{Node}. En consecuencia, cada instancia de \texttt{Rule} en el arreglo \texttt{Scene::rules} es el nodo raíz del árbol de parsing de alguna regla definida en la entrada.

Existe aproximadamente una subclase de \texttt{Node} por cada producción en la gramática. Por ejemplo, las producciones \texttt{element\_and} y \texttt{element\_or} se corresponden con las subclases \texttt{And} y \texttt{Or}, respectivamente.

La librería PLY permite asociar acciones semánticas a cada producción. Por ejemplo, la producción \texttt{arith\_expr\_number} está asociada a la siguiente función:

\begin{verbatim}
  def p_arith_expr_number(self, t):
    'arith_expr_number : NUMBER'
    t[0] = scene.Number(self.scene, t[1])
\end{verbatim}

Esta función asocia al resultado de reconocer esta producción (\texttt{t[0]}) una instancia de la clase \texttt{Number} (subclase de \texttt{Node}) que recibe el valor numérico del token \texttt{NUMBER} (\texttt{t[1]}).

De forma similar, la producción \texttt{arith\_expr\_plus} está asociada a la siguiente función:

\begin{verbatim}
  def p_arith_expr_plus(self, t):
    'arith_expr_plus : arith_expr PLUS arith_expr'
    t[0] = scene.Plus(self.scene, t[1], t[3])
\end{verbatim}

Notar que el lado derecho de esta producción incluye nodos no terminales. En este caso, los contenidos de \texttt{t[1]} y \texttt{t[3]} serán instancias de subclases de \texttt{Node} creadas por el código asociado a producciones \texttt{arith\_expr}. La función termina asignándole a \texttt{t[0]} una instancia de \texttt{Plus}, subclase de \texttt{Node}, que tiene como hijos a las dos instancias de \texttt{Node} resultantes de los nodos no terminales.

El resto de los nodos de un árbol de parsing de una regla cuaqluiera se instancian de forma similar.

A continuación se muestra una entrada de ejemplo \texttt{eg09.peg}:

\begin{verbatim}
  $ = box:tx -1 & ball:tx 1
\end{verbatim}

Y su correspondiente árbol de parsing, obtenido pasándole el parámetro \texttt{-p} al comando \texttt{pegv}:

\begin{verbatim}
Rule "$"
  And
    TX
      Box
      UMinus
        Number "1.0"
    TX
      Ball
      Number "1.0"
\end{verbatim}


\subsection{Rendering}

Pendiente.


\section{Ejemplificación de árboles de derivación}

\textbf{Entrada:}

\texttt{\$ = box : tx 1}

Árbol de derivación:

\hspace{-5cm}
\Tree [
  .rules
    [.rule\_definition
      {RULE\\(\$)}
      !\qsetw{-5cm}
      EQUALS
      !\qsetw{-5cm}
      [.element
        [.transform
          [.element
            [.primitive BOX ]
          ]
          COLON
          [.transform\_name TX ]
          [.arith\_expr
            [.arith\_expr\_number {NUMBER\\(1)} ]
          ]
        ]
      ]
    ]
    !\qsetw{5cm}
    [.rules empty ]
]

\textbf{Entrada:}

\texttt{myrule = BALL}

\texttt{\$ = myrule : s 2}

Árbol de derivación:

\hspace{-5cm}
\Tree [
  .rules
    [.rule\_definition
      {RULE\\(myrule)}
      !\qsetw{-5cm}
      EQUALS
      [.element
        BALL
      ]
    ]
    [.rules
      [.rule\_definition
        {RULE\\(\$)}
        !\qsetw{-5cm}
          EQUALS
          [.element
            [.transform
              [.element
                !\qsetw{-5cm}
                {RULE\\(myrule)}
              ]
              !\qsetw{-5cm}
              COLON
              !\qsetw{-5cm}
              [.transform\_name S ]
              [.arith\_expr
                  [.arith\_expr\_number {NUMBER\\(2)} ]
                ]
            ]
          ]
      ]
    !\qsetw{5cm}
    [.rules empty ]
    ]
]

\textbf{Entrada:}

\texttt{\$ = box | ball | \$}

Árbol de derivación:

\hspace{-5cm}
\Tree [
  .rules
    [.rule\_definition
      {RULE\\(\$)}
      !\qsetw{-5cm}
      EQUALS
      [.element
        [.element\_or
          [.element
            [.primitive BOX ]
          ]
          OR
          [.element
            [.element\_or
              [.element
                [.primitive BOX ]
              ]
              OR
              [.element
                [.rule
                  {RULE\\(\$)}
                ]
              ]
            ]
          ]
        ]
      ]
    ]
    !\qsetw{5cm}
    [.rules empty ]
]

\section{Ejemplos de programas válidos e inválidos}

\subsection{Programas válidos}

\noindent
Entrada:

\begin{verbatim}
balloon = ball:cg0:cb0:sx0.9:sz0.9                  "el globo"
        & ball:cg0:cb0:s0.04:ty-0.52                "el nudo"
        & box:cg0:cb0:s0.03:rz45:rx45:ty-0.56       "el pico"
        & line:ty-0.56:d10                          "el hilo"

line = box:sx0.01:sy0.1:sz0.01:ty-0.05 & [ line:rz10:ty-0.1 | line:rz-10:ty-0.1 ]

$ = balloon

\end{verbatim}

\grafico{ejemplo23.png}{eg23.peg}{Ej1}


\subsection{Programas inválidos}

\noindent
Entrada:
\begin{verbatim}
$ = ball : tx 2 : ts 1 : tl 2 : tz 1
\end{verbatim}

\noindent
Salida:
\begin{verbatim}
Syntax error at line 1, column 19:
$ = ball : tx 2 : ts 1 : tl 2 : tz 1
                  ^
\end{verbatim}
Los tokens ts y tl son inválidos, el parser detecta que ts no es un token y muestra la posición del token erróneo.

\noindent
Entrada:
\begin{verbatim}
ejes = box:sy0.05:sz0.05:cg0:cb0:tx0.5
$ = ejjjjes : tx 2
\end{verbatim}

\noindent
Salida:
\begin{verbatim}
LookupError: Rule ejjjjes not found.
\end{verbatim}

La regla ejjjjes no produce un error gramatical, pero si ocurre un error a la hora de realizar el renderizado, la regla no se encuntra definida, entonces se levanta un error indicando el problema.


\section{Modo de uso}

El programa admite como input, archivos con texto plano con las instrucciones para realizar el renderizado según el lenguaje.

Puede ejecutarse en la carpeta /src mediante el comando \verb+python2 main.py input_file.peg+

\section{Requerimientos}
Se requiere la utilización de Python 2.7 para la ejecución del programa.

Son necesarias además, las siguientes bibliotecas:
\begin{itemize}
\item Panda3D para el renderizado gráfico, descargable de https://www.panda3d.org/
\item Ply para el lado del parsing, descargable de http://www.dabeaz.com/ply/
\end{itemize}

El código del programa se encuentra en /src y debe ser ejecutado desde main.py

\section{Código}

\subsection{Main}

\begin{small}
\verbatiminput{../src/main.py}
\end{small}


\subsection{Lexer}

\begin{small}
\verbatiminput{../src/lexer.py}
\end{small}

\subsection{Parser}

\begin{small}
\verbatiminput{../src/parser.py}
\end{small}

\subsection{Scene}

\begin{small}
\verbatiminput{../src/scene.py}
\end{small}


\end{document}
